\documentclass[11pt]{article}

\usepackage[a4paper,margin=0.72in]{geometry}
\usepackage{amsmath,amssymb}
\usepackage{booktabs}
\usepackage{tabularx}
\usepackage{array}
\usepackage{enumitem}
\usepackage{xcolor}
\usepackage{microtype}
\usepackage{fancyhdr}
\usepackage{listings}

% -----------------------------------------------------------------------------
% Answer toggle
% Set \showanswersfalse for a student/question-only copy.
% Set \showanswerstrue  for the TA answer-key copy.
% -----------------------------------------------------------------------------
\newif\ifshowanswers
\showanswerstrue

\pagestyle{fancy}
\fancyhf{}
\lhead{DA3410}
\rhead{Viva Question Bank}
\cfoot{\thepage}

\setlength{\parindent}{0pt}
\setlength{\parskip}{4pt}
\setlist[itemize]{leftmargin=1.4em,itemsep=1pt,topsep=2pt}
\setlist[enumerate]{leftmargin=1.6em,itemsep=1pt,topsep=2pt}

\lstset{
  language=Python,
  basicstyle=\ttfamily\small,
  columns=fullflexible,
  keepspaces=true,
  showstringspaces=false,
  frame=single,
  framerule=0.3pt,
  aboveskip=4pt,
  belowskip=4pt,
  xleftmargin=0.5em,
  xrightmargin=0.5em
}

\newenvironment{question}[3]{%
  \par\medskip
  \noindent\textbf{Q#1. #2}\hfill\textit{#3}\par\smallskip
}{\par\medskip}

\newcommand{\answer}[1]{%
  \ifshowanswers
    \begin{quote}\small\textbf{Expected answer:} #1\end{quote}
  \fi
}

\newcommand{\marking}[1]{%
  \ifshowanswers
    \begin{quote}\small\textbf{1 mark:} #1\end{quote}
  \fi
}

\newcommand{\casesnote}{\textit{TA: choose one case only.}}
\newcommand{\oralnote}{\textit{Student answers orally; no code writing required.}}

\begin{document}

\begin{center}
    {\LARGE \textbf{DA3410 Viva Question Bank}}\\[4pt]
\end{center}

\section*{A. Representations, Preprocessing, and Evaluation}

\begin{question}{1}{Token frequencies and order}{Core}
A normalized character-frequency representation is built for each sequence. Can it distinguish the two sequences in the chosen case? \casesnote

\begin{center}
\begin{tabular}{cll}
\toprule
Case & Sequence 1 & Sequence 2 \\
\midrule
A & \texttt{AAAACCCC} & \texttt{ACACACAC} \\
B & \texttt{AABB} & \texttt{ABAB} \\
C & \texttt{ABCABC} & \texttt{CBAABC} \\
\bottomrule
\end{tabular}
\end{center}

\answer{No in all three cases. The sequences have the same character counts, so normalized token frequencies are identical even though their order differs.}
\marking{Student states that token-frequency features are order-invariant and therefore cannot distinguish the pair.}
\end{question}

\begin{question}{2}{Extracting $k$-mers}{Core}
List the contiguous $k$-mers for the chosen case. \casesnote

\begin{center}
\begin{tabular}{ccc}
\toprule
Case & Sequence & $k$ \\
\midrule
A & \texttt{ABCD} & 2 \\
B & \texttt{AABA} & 2 \\
C & \texttt{ABCDE} & 3 \\
\bottomrule
\end{tabular}
\end{center}

\answer{A: \texttt{AB, BC, CD}. B: \texttt{AA, AB, BA}. C: \texttt{ABC, BCD, CDE}.}
\marking{All contiguous $k$-mers for the chosen case are listed correctly.}
\end{question}

\begin{question}{3}{Why use $\log_{10}(k_{cat})$?}{Intermediate}
The raw $k_{cat}$ target spans many orders of magnitude. Why is training directly with squared error on the raw target problematic, and why does the log target help?

\answer{Very large raw values dominate squared error. Taking $\log_{10}$ compresses the range so extreme values do not dominate the loss to the same extent.}
\marking{Student connects orders-of-magnitude variation to domination of raw squared error and explains the role of the log transform.}
\end{question}

\begin{question}{4}{Train-only preprocessing}{Core}
A student wants to fit a scaler or PCA before evaluating on validation data. Which data should be used to fit the scaler/PCA?

\answer{Training data only; the fitted transform is then applied unchanged to validation and test data.}
\marking{Student says training split only.}
\end{question}

\begin{question}{5}{ROC-AUC and thresholds}{Intermediate}
For the image classification task, why can ROC-AUC be reported without choosing a classification threshold, while sensitivity and specificity depend on a threshold?

\answer{ROC-AUC evaluates the ranking of scores across thresholds. Sensitivity and specificity are computed after converting scores to class decisions at a chosen threshold.}
\marking{Student distinguishes threshold-independent ranking evaluation from threshold-dependent class decisions.}
\end{question}

\begin{question}{6}{Preprocessing-order bug}{Intermediate}
There is one evaluation bug in the code. Identify it. \oralnote

\begin{lstlisting}
X = scaler.fit_transform(X)
X_train, X_val, y_train, y_val = train_test_split(X, y)
model.fit(X_train, y_train)
\end{lstlisting}

\answer{The scaler is fitted before the train/validation split, so validation information leaks into preprocessing. Split first, fit the scaler on $X_{train}$ only, then transform validation data using that fitted scaler.}
\marking{Student identifies train--validation leakage caused by fitting preprocessing before the split.}
\end{question}

\begin{question}{7}{Reduction: mean versus sum}{Intermediate}
What changes if \texttt{mean()} is replaced by \texttt{sum()} in this loss? \oralnote

\begin{lstlisting}
per_sample = (pred - target) ** 2
loss = per_sample.mean()
\end{lstlisting}

\answer{The loss and gradients become proportional to the batch size instead of being averaged per sample. With \texttt{sum()}, larger batches generally produce larger gradient magnitudes.}
\marking{Student explains that summation changes gradient scale with batch size; the required comparison uses a mean loss.}
\end{question}

\section*{B. Activation Functions and Tensor Operations}

\begin{question}{8}{ReLU and Leaky ReLU}{Core}
Give the activation output and local derivative for the chosen case. Use the convention that the ReLU derivative at zero is $0$. \casesnote

\begin{center}
\begin{tabular}{clc}
\toprule
Case & Activation and input & Required output \\
\midrule
A & ReLU, $x=-2$ & $f(x)$ and $f'(x)$ \\
B & ReLU, $x=3$ & $f(x)$ and $f'(x)$ \\
C & ReLU, $x=0$ & $f(x)$ and $f'(x)$ \\
D & Leaky ReLU, $x=-2$, $\alpha=0.1$ & $f(x)$ and $f'(x)$ \\
\bottomrule
\end{tabular}
\end{center}

\answer{A: $(0,0)$. B: $(3,1)$. C: $(0,0)$. D: $(-0.2,0.1)$.}
\marking{Correct activation and derivative for the chosen case.}
\end{question}

\begin{question}{9}{Softmax output}{Core}
Consider the following code:

\begin{verbatim}
logits = torch.tensor(CASE, dtype=torch.float32)
p = torch.softmax(logits, dim=0)
print(p)
\end{verbatim}

What does \texttt{p} contain for the chosen case? \casesnote

\begin{center}
\begin{tabular}{cl}
\toprule
Case & \texttt{CASE} \\
\midrule
A & \texttt{[0., 0.]} \\
B & \texttt{[2., 2., 2.]} \\
C & \texttt{[-3., -3., -3., -3.]} \\
\bottomrule
\end{tabular}
\end{center}

\answer{
A: approximately \texttt{[0.5, 0.5]}. \\
B: approximately \texttt{[0.333, 0.333, 0.333]}. \\
C: approximately \texttt{[0.25, 0.25, 0.25, 0.25]}. \\
When all logits are equal, softmax assigns equal probability to every class.
}
\marking{Correct output for the chosen case.}
\end{question}

\begin{question}{10}{Stacking linear layers}{Intermediate}
The TA chooses one of the following implementations.

\textbf{Case A}
\begin{verbatim}
h = x @ W1 + b1
y = h @ W2 + b2
\end{verbatim}

\textbf{Case B}
\begin{verbatim}
h = torch.relu(x @ W1 + b1)
y = h @ W2 + b2
\end{verbatim}

For the chosen case, can the complete mapping from \texttt{x} to \texttt{y}
always be replaced by a single affine layer of the form

\begin{verbatim}
y = x @ W + b
\end{verbatim}

without changing the function? Explain briefly.

\answer{
Case A: Yes. The two affine transformations can be combined into one affine
transformation.

Case B: No, in general. ReLU introduces a nonlinear operation between the
layers, so the overall function cannot generally be collapsed into one affine
layer.
}
\marking{Correct yes/no answer for the chosen case and identifies the presence
or absence of the nonlinear activation as the reason.}
\end{question}

\begin{question}{11}{Sigmoid implementation bug}{Core}
There is one bug in this implementation. What is it? \oralnote

\begin{lstlisting}
def sigmoid(x):
    return 1 / 1 + torch.exp(-x)
\end{lstlisting}

\answer{The denominator needs parentheses: mathematically it should be $1/(1+e^{-x})$. As written, Python evaluates \texttt{1 / 1} first and then adds \texttt{exp(-x)}.}
\marking{Student identifies the missing parentheses/operator-precedence error.}
\end{question}

\begin{question}{12}{Stable softmax}{Intermediate}
The second line is replaced by \texttt{$exp_z$ = torch.exp(z)}. What changes? \oralnote

\begin{lstlisting}
z = torch.tensor([1000.0, 1001.0])
exp_z = torch.exp(z - z.max())
p = exp_z / exp_z.sum()
\end{lstlisting}

\answer{In exact mathematics the probabilities are unchanged by subtracting the same maximum from all logits, but removing the subtraction can cause overflow and non-finite values.}
\marking{Student explains numerical stability and that subtracting the maximum does not change the mathematical softmax probabilities.}
\end{question}

\begin{question}{13}{Argmax dimension}{Intermediate}
What shape does \texttt{pred} have? What would change if \texttt{dim=1} were replaced by \texttt{dim=0}? \oralnote

\begin{lstlisting}
logits = torch.randn(8, 7)
pred = logits.argmax(dim=1)
\end{lstlisting}

\answer{With \texttt{dim=1}, \texttt{pred} has shape $(8,)$: one predicted class per sample. With \texttt{dim=0}, the result has shape $(7,)$ and chooses one row/sample index for each class column, which is not the intended prediction operation.}
\marking{Student gives the correct shapes and understands which axis corresponds to classes.}
\end{question}

\section*{C. Assignment 3: Width, Depth, Capacity, and Model Selection}

\begin{question}{14}{Checkpointing implementation}{Intermediate}
The TA chooses one case. Is the checkpointing code correct? If not, identify the
problem and state the fix conceptually. \casesnote

\textbf{Case A}
\begin{verbatim}
best_loss = float("inf")

for epoch in range(epochs):
    ...
    if val_loss < best_loss:
        best_loss = val_loss
        best_params = params
\end{verbatim}

\textbf{Case B}
\begin{verbatim}
for epoch in range(epochs):
    best_loss = float("inf")
    ...
    if val_loss < best_loss:
        best_loss = val_loss
        best_params = clone_params(params)
\end{verbatim}

\textbf{Case C}
\begin{verbatim}
best_loss = float("inf")

for epoch in range(epochs):
    ...
    if train_loss < best_loss:
        best_loss = train_loss
        best_params = clone_params(params)
\end{verbatim}

\answer{
Case A: Incorrect if \texttt{best\_params = params} only stores references to
the live tensors. Later updates may also change what is considered the saved
checkpoint. Store independent copies/clones of the parameter values.

Case B: Incorrect. \texttt{best\_loss} is reset at the start of every epoch, so
almost every epoch can overwrite the previous checkpoint. Initialize it once
before the training loop.

Case C: Incorrect for the required checkpointing rule. The model is being
selected using training loss rather than validation loss. Compare
\texttt{val\_loss} against the best validation loss.
}

\marking{Student correctly identifies the error in the chosen case and gives
the conceptual correction.}
\end{question}

\begin{question}{15}{Validation-loop bugs}{Intermediate}
The TA chooses one case. Identify the problem in the code and explain what is
wrong with the resulting validation measurement. \casesnote

\textbf{Case A}
\begin{verbatim}
for xb, yb in val_loader:
    pred = forward(xb, params)
    loss = loss_fn(pred, yb)
    backward(loss)
    update(params, lr)
\end{verbatim}

\textbf{Case B}
\begin{verbatim}
for xb, yb in val_loader:
    pred = forward(xb, params)
    loss = loss_fn(pred, yb)

val_loss = train_loss
\end{verbatim}

\textbf{Case C}
\begin{verbatim}
for xb, yb in train_loader:
    train_step(xb, yb)

val_loss = evaluate(train_loader, params)
\end{verbatim}

\answer{
Case A: The model is being updated using validation samples. Validation data
must only be used for evaluation, not parameter updates. This leaks validation
information into training.

Case B: The validation predictions may be computed, but the reported
\texttt{val\_loss} is actually the training loss. The validation losses must be
accumulated/averaged separately.

Case C: The so-called validation loss is being evaluated on the training
loader. It therefore measures training performance, not generalization to the
validation set.
}

\marking{Student identifies the error in the chosen case and explains why the
reported validation performance is invalid or misleading.}
\end{question}

\begin{question}{16}{Linear-layer output shape}{Core}
What is the shape of \texttt{z}? \casesnote \ \oralnote

\begin{lstlisting}
x = torch.randn(B, D)
W = torch.randn(D, H)
b = torch.randn(H)
z = x @ W + b
\end{lstlisting}

\begin{center}
\begin{tabular}{cccc}
\toprule
Case & $B$ & $D$ & $H$ \\
\midrule
A & 4 & 3 & 5 \\
B & 8 & 10 & 2 \\
C & 1 & 6 & 7 \\
\bottomrule
\end{tabular}
\end{center}

\answer{The shape is always $(B,H)$. A: $(4,5)$. B: $(8,2)$. C: $(1,7)$.}
\marking{Correct output shape for the chosen case.}
\end{question}

\section*{D. Assignment 4: Manual Training, Autodiff, Optimization, and Checkpointing}

\begin{question}{18}{What must be saved for backward?}{Intermediate}
A manual autodiff system performs the forward operation shown below. Assume the
computation graph already records the operation and its parent tensors.

For the chosen case, is the saved information sufficient to compute the local
backward pass? If not, state what additional quantity must be retained.
\casesnote

\textbf{Case A}
\begin{verbatim}
y = a * b
saved = {"a": a}
\end{verbatim}

\textbf{Case B}
\begin{verbatim}
h = torch.maximum(z, torch.tensor(0.0))
saved = {"mask": z > 0}
\end{verbatim}

\textbf{Case C}
\begin{verbatim}
s = torch.sigmoid(z)
saved = {"s": s}
\end{verbatim}

\answer{
Case A: Not sufficient. To compute both local gradients,
$\partial y/\partial a=b$ and $\partial y/\partial b=a$, the backward operation
also needs \texttt{b} (or equivalent retained information).

Case B: Sufficient. The ReLU mask contains the information needed to multiply
the upstream gradient by $1$ where $z>0$ and $0$ elsewhere.

Case C: Sufficient. The sigmoid output is enough because its derivative can be
computed as $s(1-s)$.
}

\marking{Student correctly determines whether the saved information is
sufficient and, when necessary, identifies the missing quantity.}
\end{question}

\begin{question}{19}{Backward traversal of the computation graph}{Intermediate}
Consider the following manual autodiff forward pass:

\begin{verbatim}
a = multiply(x, w)
b = relu(a)
loss = square(b)
\end{verbatim}

The operations are recorded during the forward pass as

\begin{verbatim}
ops = [multiply_op, relu_op, square_op]
\end{verbatim}

The TA chooses one proposed backward implementation. Is it correct? If not,
state the required change. \casesnote

\textbf{Case A}
\begin{verbatim}
for op in ops:
    op.backward()
\end{verbatim}

\textbf{Case B}
\begin{verbatim}
for op in reversed(ops):
    op.backward()
\end{verbatim}

\textbf{Case C}
\begin{verbatim}
for op in reversed(ops):
    op.backward()
    op.parents[0].grad = op.grad
\end{verbatim}

\answer{
Case A: Incorrect. The operations are being differentiated in forward order.
Backward must begin at the loss and traverse the recorded operations in reverse
order.

Case B: Correct, assuming each operation's backward function propagates and
accumulates its gradients into its parents.

Case C: Potentially incorrect. A parent's gradient should generally be
accumulated rather than blindly overwritten, because the same tensor may
contribute to the loss through multiple paths in the computation graph.
}

\marking{Student correctly judges the chosen implementation and identifies
reverse traversal and, for Case C, gradient accumulation as the key issue.}
\end{question}

\begin{question}{20}{One-operation backward: multiplication}{Core}
For the chosen case, give $\partial L/\partial a$ and $\partial L/\partial b$. \casesnote \ \oralnote

\begin{lstlisting}
y = a * b
# upstream gradient is g = dL/dy
\end{lstlisting}

\begin{center}
\begin{tabular}{cccc}
\toprule
Case & $a$ & $b$ & $g$ \\
\midrule
A & 3 & 2 & 1 \\
B & 4 & 5 & 2 \\
C & -2 & 3 & 1 \\
\bottomrule
\end{tabular}
\end{center}

\answer{Use $\partial L/\partial a=g\,b$ and $\partial L/\partial b=g\,a$. A: $(2,3)$. B: $(10,8)$. C: $(3,-2)$.}
\marking{Correct two gradients for the chosen case.}
\end{question}

\begin{question}{21}{Gradient accumulation at a branch}{Intermediate}
There is one bug in the backward code. Identify it. \oralnote

\begin{lstlisting}
a = x * x
L = a + x

grad_x = 1.0          # path L -> x
grad_x = 2 * x        # path L -> a -> x
\end{lstlisting}

\answer{The second assignment overwrites the gradient from the first path. Because $x$ influences $L$ through two paths, the contributions must be added: $dL/dx = 1 + 2x$.}
\marking{Student identifies that gradients from multiple graph paths must accumulate rather than overwrite one another.}
\end{question}

\begin{question}{22}{Linear backward: weight-gradient shape}{Intermediate}
What is the shape of \texttt{grad\_W}? \oralnote

\begin{lstlisting}
x = torch.randn(16, 10)
grad_z = torch.randn(16, 7)
grad_W = x.T @ grad_z
\end{lstlisting}

\answer{\texttt{x.T} has shape $(10,16)$ and \texttt{grad\_z} has shape $(16,7)$, so \texttt{grad\_W} has shape $(10,7)$, matching the weight matrix shape.}
\marking{Student gives $(10,7)$.}
\end{question}

\begin{question}{23}{Linear backward: bias reduction}{Intermediate}
Why is the bias gradient reduced with \texttt{sum(dim=0)} rather than \texttt{sum(dim=1)}? What shape does it have? \oralnote

\begin{lstlisting}
grad_z = torch.randn(32, 128)
grad_b = grad_z.sum(dim=0)
\end{lstlisting}

\answer{Each of the 128 bias entries is shared across all 32 samples, so gradient contributions must be summed over the batch dimension. The resulting shape is $(128,)$.}
\marking{Student identifies the batch dimension and gives shape $(128,)$.}
\end{question}

\begin{question}{24}{Which input-gradient implementation is correct?}{Core}
For \texttt{z = x @ W}, which candidate correctly computes the gradient with respect to \texttt{x}? \oralnote

\begin{lstlisting}
# Candidate A
 grad_x = grad_z @ W.T

# Candidate B
 grad_x = W.T @ grad_z

# Candidate C
 grad_x = grad_z.T @ W
\end{lstlisting}

\answer{Candidate A. If $x$ is $(B,D)$, $W$ is $(D,H)$, and $grad\_z$ is $(B,H)$, then $grad\_z W^T$ has shape $(B,D)$, matching $x$.}
\marking{Student chooses A and can justify it by the operation/shape.}
\end{question}

\begin{question}{25}{Updates and mini-batch construction}{Intermediate}
Consider the following code for one training epoch:

\begin{verbatim}
updates = 0

for start in range(0, N, batch_size):
    xb = X[start:start + batch_size]
    update(params, xb)
    updates += 1

print(updates, len(xb))
\end{verbatim}

For the chosen case, what does the final \texttt{print} output? \casesnote

\begin{center}
\begin{tabular}{ccc}
\toprule
Case & \texttt{N} & \texttt{batch\_size} \\
\midrule
A & 128 & 1 \\
B & 256 & 128 \\
C & 300 & 128 \\
D & 300 & 300 \\
\bottomrule
\end{tabular}
\end{center}

\answer{
A: \texttt{128 1}. There are 128 SGD updates. \\
B: \texttt{2 128}. There are two complete mini-batches. \\
C: \texttt{3 44}. The first two batches contain 128 samples and the final
batch contains 44. \\
D: \texttt{1 300}. This is one full-batch update.
}

\marking{Student gives both the correct number of updates and the size of the
final batch for the chosen case.}
\end{question}

\begin{question}{26}{Mini-batch construction}{Core}
For the chosen case, how many batches are produced and what is the size of the last batch? \casesnote \ \oralnote

\begin{lstlisting}
for i in range(0, N, batch_size):
    xb = X[i:i + batch_size]
\end{lstlisting}

\begin{center}
\begin{tabular}{ccc}
\toprule
Case & $N$ & \texttt{batch\_size} \\
\midrule
A & 10 & 4 \\
B & 8 & 4 \\
C & 9 & 8 \\
\bottomrule
\end{tabular}
\end{center}

\answer{A: 3 batches, last size 2. B: 2 batches, last size 4. C: 2 batches, last size 1.}
\marking{Correct number of batches and final-batch size.}
\end{question}

\begin{question}{27}{Shuffling bug}{Intermediate}
The code creates a random permutation, but the data are still processed in their original order. Where is the bug? \oralnote

\begin{lstlisting}
perm = torch.randperm(len(X))
for i in range(0, len(X), batch_size):
    xb = X[i:i + batch_size]
    yb = y[i:i + batch_size]
\end{lstlisting}

\answer{\texttt{perm} is never used. The batches must index the shuffled sample indices, for example by first applying the permutation to $X$ and $y$, or by using the appropriate slice of \texttt{perm} for indexing.}
\marking{Student identifies that generating a permutation alone does not shuffle the training samples.}
\end{question}

\begin{question}{28}{Noise across optimization regimes}{Core}
Rank SGD, mini-batch gradient descent, and full-batch gradient descent from most update noise to least update noise.

\answer{SGD is most noisy, mini-batch is intermediate, and full-batch is least noisy because each update uses the complete training set.}
\marking{Correct ordering.}
\end{question}

\begin{question}{29}{Checkpoint aliasing bug}{Intermediate}
The validation loss improves, but the supposedly saved ``best'' parameters keep changing during later training. Identify the bug. \oralnote

\begin{lstlisting}
if val_loss < best_loss:
    best_loss = val_loss
    best_params = params

# later: parameters are updated in place
\end{lstlisting}

\answer{\texttt{best\_params} refers to the same mutable parameter objects as \texttt{params}, so later in-place updates also change the saved checkpoint. The tensors must be copied/cloned when the checkpoint is stored.}
\marking{Student identifies reference/aliasing of mutable tensors and the need to store independent copies.}
\end{question}

\begin{question}{30}{Debugging an incorrect backward update}{Intermediate}
A student computes the gradient for the last layer, immediately updates that layer's weights, and then uses the \emph{updated} weights to propagate gradients into earlier layers. What is wrong with this procedure?

\answer{The backward pass is no longer computing gradients for the same parameter state used in the forward pass. Gradients should be computed first, then parameters should be updated.}
\marking{Student identifies that all gradients must correspond to the original forward-pass parameters before any update is applied.}
\end{question}

\section*{E. Gradient Diagnostics and Generalization}

\begin{question}{31}{Gradient-propagation diagnostic code}{Intermediate}
The following code is used to initialize a deep network before measuring
layerwise gradient magnitudes.

For the chosen case, state whether the setup is intended to produce
\textbf{stable}, \textbf{vanishing}, or \textbf{exploding} gradient behaviour.
\casesnote

\textbf{Case A}
\begin{verbatim}
for l in range(depth):
    W[l] = he_init(fan_in, fan_out)
activation = relu
\end{verbatim}

\textbf{Case B}
\begin{verbatim}
for l in range(depth):
    W[l] = xavier_init(fan_in, fan_out)
activation = sigmoid
\end{verbatim}

\textbf{Case C}
\begin{verbatim}
for l in range(depth):
    W[l] = 2 * he_init(fan_in, fan_out)
activation = relu
\end{verbatim}

\answer{
Case A: Stable reference. ReLU with He initialization is intended to maintain
reasonable activation and gradient magnitudes through depth.

Case B: Vanishing-gradient condition. Repeated sigmoid derivatives can reduce
the gradient as it propagates through many layers.

Case C: Exploding-gradient condition. Doubling the He-initialized weights
amplifies signals repeatedly through the deep network and is intended to
produce growing gradients.
}

\marking{Student correctly identifies the intended gradient behaviour for the
chosen code and gives the relevant reason.}
\end{question}

\begin{question}{32}{Group leakage in train--validation splitting}{Intermediate}
Suppose \texttt{group\_id[i]} identifies which related group sample
\texttt{i} belongs to. Samples with the same \texttt{group\_id} should not
appear in both training and validation.

For the chosen implementation, does the code guarantee this? If not, explain
the problem. \casesnote

\textbf{Case A}
\begin{verbatim}
perm = torch.randperm(len(X))
cut = int(0.8 * len(X))

train_idx = perm[:cut]
val_idx   = perm[cut:]
\end{verbatim}

\textbf{Case B}
\begin{verbatim}
groups = torch.unique(group_id)
groups = groups[torch.randperm(len(groups))]
cut = int(0.8 * len(groups))

train_groups = groups[:cut]
val_groups   = groups[cut:]
\end{verbatim}

The corresponding samples are then selected according to their group.

\textbf{Case C}
\begin{verbatim}
train_idx = torch.where(group_id != 5)[0]
val_idx   = torch.where(group_id == 5)[0]
\end{verbatim}

\answer{
Case A: No. It randomly splits individual samples, so samples belonging to the
same group may appear in both training and validation. This can give an
overly optimistic validation result.

Case B: Yes. The split is performed on unique groups first, so an entire group
is assigned to only one split.

Case C: There is no group overlap between the two sets, so it avoids that
specific leakage. However, it uses only one group for validation and therefore
may be a poor or unrepresentative validation split.
}

\marking{Student correctly judges whether the chosen implementation separates
related groups and explains the consequence of any problem.}
\end{question}

\end{document}
